% Options for packages loaded elsewhere
\PassOptionsToPackage{unicode}{hyperref}
\PassOptionsToPackage{hyphens}{url}
\documentclass[
  a4paper,
]{article}
\usepackage{xcolor}
\usepackage[margin=25mm]{geometry}
\usepackage{amsmath,amssymb}
\setcounter{secnumdepth}{-\maxdimen} % remove section numbering
\usepackage{iftex}
\ifPDFTeX
  \usepackage[T1]{fontenc}
  \usepackage[utf8]{inputenc}
  \usepackage{textcomp} % provide euro and other symbols
\else % if luatex or xetex
  \usepackage{unicode-math} % this also loads fontspec
  \defaultfontfeatures{Scale=MatchLowercase}
  \defaultfontfeatures[\rmfamily]{Ligatures=TeX,Scale=1}
\fi
\usepackage{lmodern}
\ifPDFTeX\else
  % xetex/luatex font selection
\fi
% Use upquote if available, for straight quotes in verbatim environments
\IfFileExists{upquote.sty}{\usepackage{upquote}}{}
\IfFileExists{microtype.sty}{% use microtype if available
  \usepackage[]{microtype}
  \UseMicrotypeSet[protrusion]{basicmath} % disable protrusion for tt fonts
}{}
\makeatletter
\@ifundefined{KOMAClassName}{% if non-KOMA class
  \IfFileExists{parskip.sty}{%
    \usepackage{parskip}
  }{% else
    \setlength{\parindent}{0pt}
    \setlength{\parskip}{6pt plus 2pt minus 1pt}}
}{% if KOMA class
  \KOMAoptions{parskip=half}}
\makeatother
\usepackage{longtable,booktabs,array}
\newcounter{none} % for unnumbered tables
\usepackage{calc} % for calculating minipage widths
% Correct order of tables after \paragraph or \subparagraph
\usepackage{etoolbox}
\makeatletter
\patchcmd\longtable{\par}{\if@noskipsec\mbox{}\fi\par}{}{}
\makeatother
% Allow footnotes in longtable head/foot
\IfFileExists{footnotehyper.sty}{\usepackage{footnotehyper}}{\usepackage{footnote}}
\makesavenoteenv{longtable}
\usepackage{graphicx}
\makeatletter
\newsavebox\pandoc@box
\newcommand*\pandocbounded[1]{% scales image to fit in text height/width
  \sbox\pandoc@box{#1}%
  \Gscale@div\@tempa{\textheight}{\dimexpr\ht\pandoc@box+\dp\pandoc@box\relax}%
  \Gscale@div\@tempb{\linewidth}{\wd\pandoc@box}%
  \ifdim\@tempb\p@<\@tempa\p@\let\@tempa\@tempb\fi% select the smaller of both
  \ifdim\@tempa\p@<\p@\scalebox{\@tempa}{\usebox\pandoc@box}%
  \else\usebox{\pandoc@box}%
  \fi%
}
% Set default figure placement to htbp
\def\fps@figure{htbp}
\makeatother
\setlength{\emergencystretch}{3em} % prevent overfull lines
\providecommand{\tightlist}{%
  \setlength{\itemsep}{0pt}\setlength{\parskip}{0pt}}
\usepackage{bookmark}
\IfFileExists{xurl.sty}{\usepackage{xurl}}{} % add URL line breaks if available
\urlstyle{same}
\hypersetup{
  pdftitle={Reading Relations{,} Loops{,} Flux{,} and Condensates from the Closure of Complex Numbers --- A Minimal-Assumption Thought Experiment That Shaves Off Axiom Candidates},
  pdfauthor={Noriaki Kihara (WF System Co., Ltd.) ORCID 0009-0004-6753-4020},
  hidelinks,
  pdfcreator={LaTeX via pandoc}}

\title{Reading Relations, Loops, Flux, and Condensates from the Closure
of Complex Numbers --- A Minimal-Assumption Thought Experiment That
Shaves Off Axiom Candidates}
\author{Noriaki Kihara (WF System Co., Ltd.) ORCID 0009-0004-6753-4020}
\date{September 19, 2026 Version DOI 10.5281/zenodo.22841917}

\begin{document}
\maketitle

\textbf{Author:} Noriaki Kihara\\
\textbf{ORCID:} 0009-0004-6753-4020\\
\textbf{Date:} 2026-09-19\\
\textbf{Version:} v1.3 (review revision: the vanishing global flux
recast as a derived identity, domain conditions for the complex
logarithm, the curvature readout restricted to a correspondence map,
finite sequences made explicit)\\
\textbf{Version DOI:} 10.5281/zenodo.22841917\\
\textbf{Concept DOI:} 10.5281/zenodo.22841916

\begin{center}\rule{0.5\linewidth}{0.5pt}\end{center}

\subsection{Abstract}\label{abstract}

This paper proposes no new algebraic law and no new equation of motion.
The starting point is only the ordinary complex number
\(z\in\mathbb C\), and the only facts used are the known ones that
complex numbers are closed under addition, multiplication, and squaring.
In particular, for any finite complex sequence \(\{z_n\}_{n=1}^{N}\) we
have \(z_n^2\in\mathbb C\), so setting \(Z_n:=z_n^2\) we can read

\[
\sum_{n=1}^{N} z_n^2=\sum_{n=1}^{N} Z_n=K,\qquad K\in\mathbb C.
\]

Furthermore, any \(K\in\mathbb C\) can be written as \(K=C^2\) for some
\(C\in\mathbb C\), so

\[
\sum_n z_n^2=C^2
\]

is not a new physical axiom but a notation that stays within the closure
of the complex numbers.

The subject of this paper is not this trivial fact itself but its
\textbf{order of readout}. In the thought experiment, we questioned, one
by one, complete networks, external standards, ratios, equality tests,
zero closure, and the presupposition of triangles and vertices, removing
hidden assumptions against namelessness. As a result, when the squared
quantities \(Z_n\) are read as ordinary complex relations, we found that
sums of multiple relations, loops, loop residuals, face fluxes,
cancellation of shared parts, closed polyhedra as condensates, and
further the hierarchization of condensates can all be described while
staying within the same complex-number type.

If the sum of three relations is zero it can be read as a loop; if
nonzero it carries a residual \(K_f\). When multiple loops are glued
together with orientations, shared edges cancel with opposite signs, and
only the quantities remaining on the external boundary survive as
effective flux. Since the face residual \(K_f\) is itself a complex
number, it can be reused as a relation at the next level of the
hierarchy. This hierarchy has the characteristic that edge, face, and
cell types need not be postulated separately: the recursion runs inside
the same \(\mathbb C\).

At the end, the paper compares structural similarities with lattice
gauge theory, coarse-graining of spin networks, closure defects, the
Bianchi identity, and bubble networks. However, none of these is placed
at the starting point. The candidate novelty of this paper is not a
simplification of existing theories but \textbf{the fact that, starting
only from the trivial closure of ordinary complex numbers and shaving
off hidden assumptions along the way, we obtain a minimal interpretive
framework that reads off loops, curvature, flux, and condensates}.

\textbf{Keywords:} complex numbers, squared quantities, loop, closure,
closure defect, flux, curvature, condensate, namelessness, minimal
assumptions, lattice gauge theory

\begin{center}\rule{0.5\linewidth}{0.5pt}\end{center}

\subsection{1. What This Paper Does and Does Not
Claim}\label{what-this-paper-does-and-does-not-claim}

\subsubsection{1.1 What is not claimed}\label{what-is-not-claimed}

This paper does not claim the following as new mathematical theorems.

\begin{enumerate}
\def\labelenumi{\arabic{enumi}.}
\tightlist
\item
  That the square of a complex number is a complex number.
\item
  That a finite sum of complex numbers is a complex number.
\item
  That every complex number has a complex square root.
\item
  That if a sum of complex vectors is zero it can be drawn as a closed
  polygon.
\item
  That on an oriented boundary a shared element appears with opposite
  signs.
\end{enumerate}

These are known mathematics.

\subsubsection{1.2 The subject of this
paper}\label{the-subject-of-this-paper}

What this paper addresses is \textbf{in what order, and without which
additional assumptions, the known facts above can be read as physical
structure}.

Therefore this paper is not a ``paper proposing a new equation'' but

\[
\boxed{\text{a minimal-assumption readout paper on known identity structures.}}
\]

\subsubsection{1.3 Methodology}\label{methodology}

At each stage of the thought experiment, before introducing a new
structure we ask:

\begin{itemize}
\tightlist
\item
  Is the structure derived from the algebra of complex numbers?
\item
  Are we implicitly importing an external standard?
\item
  Does it require names of individuals or vertices?
\item
  Are we granting in advance a complete graph, distance, dimension,
  coordinates, or time?
\item
  Are we placing a ``same/different'' detector as a primitive operation?
\item
  Are we importing an existing physical theory back into the
  foundations?
\end{itemize}

Only what cannot be derived yet is needed remains as an \textbf{axiom
candidate}; whatever can be derived is removed from the axiom
candidates.

\subsubsection{1.4 Revision note for v1.2: a hidden assumption about the
completeness of many-body
relations}\label{revision-note-for-v1.2-a-hidden-assumption-about-the-completeness-of-many-body-relations}

After completing v1.1 of this paper, we read Anegawa, Miyata, Suzuki,
Tamaoka (2026) {[}6{]}. Although that paper works in a different
framework --- quantum information and holography --- it exhibits
examples in which, for a fixed tripartition, quantities measuring
genuine multi-partite entanglement can continue to evolve in time even
after the ordinary entanglement entropies of the subsystems have
saturated.

Prompted by examining that paper, we noticed that v1.1 of the present
paper implicitly retained the assumption

\[
\boxed{\{Z_e,K_f\}\ \text{suffices to exhaust the effective relational information of a many-body system.}}
\]

That is, we had effectively presupposed that the pairwise complex
relations \(Z_e\), together with the set of loop residuals \(K_f\)
constructed from them, suffice to read the relational structure of a
many-body system. This completeness, however, has not been derived from
the complex-number algebra of this paper.

The results of Anegawa et al.~{[}6{]} neither directly support nor
refute the present model, and \(K_f\) cannot be identified with
multi-entropy. On the other hand, the concrete example that
\textbf{irreducible many-body relations can remain even after low-order
readouts saturate} shows that the above completeness cannot be taken for
granted.

Therefore, in v1.2 this assumption is removed from the derived
properties, and its validity is stated explicitly as a problem for
future work. This revision is the result of applying to the paper itself
the very methodology the paper adopts: ``when a hidden assumption is
discovered, delete it or put it on hold.''

\begin{center}\rule{0.5\linewidth}{0.5pt}\end{center}

\subsection{2. Starting Point: A Single Complex
Number}\label{starting-point-a-single-complex-number}

Take as the starting point

\[
z\in\mathbb C.
\]

Here we assume no particles, edges, vertices, faces, space, time, gauge
group, or network.

From the closure of the complex field,

\[
z^2\in\mathbb C.
\]

This is utterly trivial. But in this paper we pursue how much structure
can be read without breaking this triviality.

Even when a single complex number is multiplied by itself,

\[
z\mapsto z^2,
\]

the result is again a single complex number. Therefore, at this stage
one cannot say ``one body has split into two'' or ``a new vertex has
appeared.'' The algebraic type remains

\[
\mathbb C\to\mathbb C.
\]

What matters here is that \textbf{the generation history and the number
of individuals cannot be recovered from the results of complex
operations alone}. Hence, to discuss structure, the case in which
multiple complex relations appear simultaneously must be considered
separately.

\begin{center}\rule{0.5\linewidth}{0.5pt}\end{center}

\subsection{3. Two Complex Numbers: Take the Product, Not a
Comparator}\label{two-complex-numbers-take-the-product-not-a-comparator}

Suppose there are two complex numbers

\[
z_1,z_2\in\mathbb C.
\]

One might first consider taking the ratio \(z_2/z_1\) as the relational
quantity. However, reading a ratio as ``larger/smaller'' or
``same/different'' risks smuggling in a comparison standard. Likewise,
placing the equality test \(z_1=z_2\) as a primitive operation is too
strong for the purposes of this paper.

The product

\[
z_1z_2,
\]

on the other hand, requires no external standard. From the closure of
complex numbers,

\[
z_1z_2\in\mathbb C.
\]

Expanding into real and imaginary parts, for

\[
z_1=a+ib,\qquad z_2=c+id
\]

we have

\[
z_1z_2=(ac-bd)+i(ad+bc).
\]

The result is again a single complex number.

Using the same value twice gives

\[
z\,z=z^2,
\]

so the self-product is a special case of the two-body product. But there
is no need to detect from outside that the values are equal and split
the structure into cases. Algebraically, one simply takes the product.

Also, since the complex numbers have no zero divisors,

\[
z_1z_2=0
\]

implies that at least one factor is zero. Hence the minimal algebraic
boundary that appears without any external scale is ``zero / nonzero.''
In this paper, however, we do not elevate this to a new physical axiom.

\begin{center}\rule{0.5\linewidth}{0.5pt}\end{center}

\subsection{4. Reading the Squared Quantity as the Basic
Quantity}\label{reading-the-squared-quantity-as-the-basic-quantity}

Here we perform the main rereading of this paper.

For each complex number define

\[
\boxed{Z_n:=z_n^2}.
\]

Since the square of a complex number is a complex number,

\[
Z_n\in\mathbb C.
\]

Then the expression

\[
\sum_n z_n^2,
\]

which often looks like a special form, is simply

\[
\boxed{\sum_n Z_n},
\]

an ordinary sum of complex numbers.

Therefore

\[
\sum_n z_n^2=C^2
\]

is also, once written as

\[
\boxed{\sum_n Z_n=K},\qquad K:=C^2\in\mathbb C,
\]

just an ordinary complex sum.

At this point there is no need to treat \(C^2\) as ``special because it
is a square.'' In the space of squared quantities, \(K=C^2\) is an
ordinary complex number.

This rereading is the core of this paper.

\begin{figure}
\centering
\pandocbounded{\includegraphics[keepaspectratio,alt={Fig. 1 From the closure of complex numbers to the loop residual}]{figs/fig1_complex_closure_loop_residual_en.png}}
\caption{Fig. 1 From the closure of complex numbers to the loop
residual}
\end{figure}

\emph{Fig. 1 From the closure of complex numbers to the loop residual.
The trivial identity \(\sum z_n^2=C^2\) is reread as the sum
\(\sum Z_n=K\) of complex relations \(Z_n:=z_n^2\). \(K=0\) represents
closure; \(K\neq0\) represents a loop residual.}

\begin{center}\rule{0.5\linewidth}{0.5pt}\end{center}

\subsection{5. An Auxiliary Two-Body Thought Experiment: When the Sum Is
Real}\label{an-auxiliary-two-body-thought-experiment-when-the-sum-is-real}

Consider an auxiliary thought experiment in which, for two complex
numbers, the only constraint imposed is

\[
z_1^2+z_2^2=R^2\in\mathbb R.
\]

Writing the squared quantities as

\[
Z_1=z_1^2,\qquad Z_2=z_2^2,
\]

this is simply

\[
\boxed{Z_1+Z_2=R^2}.
\]

Subtracting the center,

\[
W_1:=Z_1-\frac{R^2}{2},\qquad
W_2:=Z_2-\frac{R^2}{2},
\]

we get

\[
\boxed{W_1+W_2=0},
\]

that is,

\[
W_2=-W_1.
\]

In the original \(z\) representation this involves square-root branches,
but once the squared quantity \(Z\) is taken as basic, the two bodies
simplify to an antisymmetric configuration about the center \(R^2/2\) on
the real axis.

This auxiliary example shows the advantage of ``reading squared
quantities as basic quantities,'' but the paper does not require
\(R^2\in\mathbb R\) throughout.

\begin{center}\rule{0.5\linewidth}{0.5pt}\end{center}

\subsection{6. Three Bodies: The First Nontrivial
Loop}\label{three-bodies-the-first-nontrivial-loop}

For three squared quantities, let

\[
Z_1+Z_2+Z_3=R^2\in\mathbb R,
\]

and set

\[
W_k:=Z_k-\frac{R^2}{3}.
\]

Then

\[
\boxed{W_1+W_2+W_3=0}.
\]

Hence any one of them is determined by the other two:

\[
W_3=-(W_1+W_2).
\]

What matters here is that there is no need to presuppose a ``triangle.''

First there are three complex relations \(W_1,W_2,W_3\), and their sum
happens to be zero. When this algebraic condition is drawn in the
complex plane, it can be read as a closed triangle.

That is, the logical order is not

\[
\boxed{\text{assume a triangle}\Rightarrow\text{loop condition}}
\]

but

\[
\boxed{W_1+W_2+W_3=0\Rightarrow\text{readable as a triangular loop}}.
\]

This ordering becomes important in the later discussion of namelessness.

\begin{figure}
\centering
\pandocbounded{\includegraphics[keepaspectratio,alt={Fig. 2 Three bodies: one is determined by the other two}]{figs/fig2_three_body_triangular_loop_en.png}}
\caption{Fig. 2 Three bodies: one is determined by the other two}
\end{figure}

\emph{Fig. 2 The triangular loop of three bodies. From the zero sum of
three complex relations, any one of them is fixed as the reverse of the
sum of the other two, and the closed triangle is read afterwards.}

\begin{center}\rule{0.5\linewidth}{0.5pt}\end{center}

\subsection{7. What Four Bodies Expose: Bodies and Relations Are Not the
Same}\label{what-four-bodies-expose-bodies-and-relations-are-not-the-same}

With three bodies, the number of vertices (3) coincides with the number
of edges (3), so it is easy to confuse ``bodies'' with ``relations.''

With four bodies, however, if one assumes all pairwise relations, the
number of relations is

\[
M=\binom{4}{2}=6.
\]

Here one is tempted to declare ``four bodies, therefore the complete
graph \(K_4\).'' But this is not allowed by the methodology of this
paper, because complete connectivity has not been derived from the
closure of complex numbers.

In fact, even in a five-relation system consisting of the four outer
edges and one diagonal on four vertices, there exist nontrivial
solutions in which two triangular loops close simultaneously.

For example, with oriented relations,

\[
W_{12}+W_{23}-W_{13}=0
\]

\[
W_{13}+W_{34}+W_{41}=0.
\]

From the first equation,

\[
W_{13}=W_{12}+W_{23},
\]

and substituting into the second,

\[
W_{12}+W_{23}+W_{34}+W_{41}=0.
\]

Therefore,

\[
\boxed{\text{triangular loops close}\not\Rightarrow\text{complete connectivity}}.
\]

By this counterexample, the complete-network hypothesis is excluded from
the thought experiment.

\begin{figure}
\centering
\pandocbounded{\includegraphics[keepaspectratio,alt={Fig. 3 Four-vertex network comparison of triangular loops}]{figs/fig3_four_vertex_triangular_loops_en.png}}
\caption{Fig. 3 Four-vertex network comparison of triangular loops}
\end{figure}

\emph{Fig. 3 On four vertices, the complete network \(K_4\) cannot be
derived from the mere existence of triangular loops. Even with 4
vertices and 5 relations two triangular loops can close, so complete
connectivity is an additional assumption.}

\begin{center}\rule{0.5\linewidth}{0.5pt}\end{center}

\subsection{8. A Dual Turn: From Vertex-Centered to
Edge-Centered}\label{a-dual-turn-from-vertex-centered-to-edge-centered}

To read structure without first laying down a complete graph, we invert
the viewpoint.

Conventional pictures tend to think

\[
\text{vertices}\to\text{edges}\to\text{faces}.
\]

In this paper, however, the complex relations \(Z_e\) are basic, so we
read instead

\[
\boxed{\text{relations (edges)}\to\text{loops}\to\text{faces}}.
\]

If three relations \(Z_1,Z_2,Z_3\) satisfy

\[
Z_1+Z_2+Z_3=0,
\]

then those three relations can be read as one closed face.

Here we do not specify in advance ``which vertices are shared.'' The
loop structure comes first; vertices and cell structure are read
afterwards, dually.

With this ordering, there is no need to postulate complete networks or
vertex names as primitive structure.

\begin{center}\rule{0.5\linewidth}{0.5pt}\end{center}

\subsection{9. Loop Residual}\label{loop-residual}

When a loop does not close completely, we write

\[
\boxed{\sum_{e\in\partial f}\epsilon_{fe}Z_e=K_f},\qquad K_f\in\mathbb C.
\]

\(\epsilon_{fe}=\pm1\) is the sign relative to the orientation of the
loop.

If \(K_f=0\) the loop closes completely; if \(K_f\neq0\) it carries a
complex residual.

What matters is that \(K_f\) is again an ordinary complex number:

\[
\boxed{Z_e\in\mathbb C\quad\Rightarrow\quad K_f\in\mathbb C}.
\]

Hence ``edge quantities'' and ``face residuals'' do not change the
algebraic type.

\begin{center}\rule{0.5\linewidth}{0.5pt}\end{center}

\subsection{10. The Residual as a Central-Projection
Radius}\label{the-residual-as-a-central-projection-radius}

Since \(K_f\in\mathbb C\), there exists \(C_f\in\mathbb C\) such that

\[
\boxed{K_f=C_f^2}.
\]

Therefore

\[
\sum_{e\in\partial f}\epsilon_{fe}Z_e=C_f^2
\]

can be read as a complex version of a central-projection radius written
in squared quantities.

At this stage there is no need to add ``curvature'' as an entirely
separate variable. Complete closure means \(C_f=0\); non-closure means a
finite complex radius.

When one reads this as a central-projection constant-curvature geometry,
for non-closed faces (\(K_f\neq0\)) the curvature scale is the inverse
squared radius:

\[
\kappa_f\propto \frac{1}{C_f^2}=\frac{1}{K_f}.
\]

This curvature readout is a correspondence map. In particular,
restricting to the real positive branch

\[
K_f=C_f^2\in\mathbb R_{>0}
\]

and adopting the standard convention for the radius of curvature, one
recovers the real constant-curvature readout

\[
\kappa_f=\frac{1}{K_f}.
\]

The geometric meaning carried by a general complex \(K_f\), in
particular \(\operatorname{Im}K_f\neq0\), is not derived in this paper
and is left as a connection problem.

What matters in this paper is that \(K_f\) was not added from outside as
a ``curvature variable'': the loop residual itself already exists as the
quantity determining the central-projection radius.

\begin{center}\rule{0.5\linewidth}{0.5pt}\end{center}

\subsection{11. Cancellation of Shared Edges and Face
Flux}\label{cancellation-of-shared-edges-and-face-flux}

Consider gluing two triangular faces along one internal edge.

Let the residuals of faces \(f_1\) and \(f_2\) be

\[
K_{f_1}=Z_{12}+Z_{23}-Z_{13}
\]

\[
K_{f_2}=Z_{13}+Z_{34}+Z_{41}.
\]

Adding the two, the shared edge \(Z_{13}\) cancels:

\[
-Z_{13}+Z_{13}=0,
\]

leaving

\[
\boxed{K_{f_1}+K_{f_2}=Z_{12}+Z_{23}+Z_{34}+Z_{41}}.
\]

That is, internal relations disappear and only the relations touching
the external boundary remain.

Reading \(K_f\) as a face flux, we obtain the discrete Stokes-type
structure

\[
\boxed{\text{cancellation of shared internal relations}\Rightarrow\text{only the external boundary flux remains}}.
\]

Here too, Stokes' theorem was not placed at the starting point. This is
the result of following the algebraic fact that when oriented loops are
glued, the shared term appears with opposite signs.

\begin{center}\rule{0.5\linewidth}{0.5pt}\end{center}

\subsection{12. Reading a Closed Polyhedron as a
Condensate}\label{reading-a-closed-polyhedron-as-a-condensate}

Consider a structure \(\mathcal C\) in which multiple faces are glued
together and the outer shell is closed.

Each face carries a complex flux \(K_f\), and on a face shared by
adjacent cells the orientations are opposite, so the contribution of an
internal face cancels:

\[
K_f+(-K_f)=0.
\]

Therefore, in the flux readout defined in this paper, shared internal
relations cancel and the quantity remaining on the boundary is read as
the effective quantity.

However, \textbf{it is not proven in this paper that this boundary flux
exhausts all irreducible relational information of the many-body
system.} The possibility that higher-order relational quantities exist
independently, not captured by pairwise relations and loop residuals, is
explicitly left as an open problem in v1.2.

This paper proposes to call this closed polyhedral relational structure
a \textbf{condensate}:

\[
\boxed{\text{condensate}=\text{a closed relational complex whose internal relations cancel on shared faces}}.
\]

This is not a definition of ``particle = point.'' A condensate is
defined as many complex relations forming a closed face structure that
can be read from outside as a single effective object.

This interpretation is a physical readout, not a mathematical identity,
and we keep it distinguished as an axiom candidate / interpretation
candidate.

\begin{center}\rule{0.5\linewidth}{0.5pt}\end{center}

\subsection{13. Local Flux and Global
Flux}\label{local-flux-and-global-flux}

Consider the closed polyhedral condensate of this paper, that is, a
closed, oriented, 2-manifold-like face complex \(\mathcal C\) in which
every edge belongs to exactly two faces with opposite orientations. Let
\(F(\mathcal C)\) denote the set of oriented faces constituting the
condensate \(\mathcal C\). With the face flux

\[
K_f=\sum_{e\in\partial f}\epsilon_{fe}Z_e,
\]

we have

\[
\sum_{f\in F(\mathcal C)}K_f=\sum_{e}\left(\sum_{f\in F(\mathcal C)}\epsilon_{fe}\right)Z_e=0.
\]

This holds because for each shared edge the incidence signs \(+1\) and
\(-1\) each appear exactly once; it is a discrete identity of
\(\partial^2=0\) type. This structure corresponds to the commutative
version of the structure in which reference {[}4{]} reads the closure
constraint as a Bianchi identity.

Therefore, the vanishing of the global flux is not an additional
condition but a consequence of this closed oriented structure. It does
not, however, imply

\[
K_f=0
\]

for each face. Faces with \(K_f\neq0\) can exist locally.

This is important: a structure is possible in which nonzero interactions
persist inside while the whole remains closed.

Therefore,

\[
\boxed{\text{global flux zero}\neq\text{local flux zero}}.
\]

If one adopts a physical correspondence to gauge theory, the nonzero
fluxes on internal shared faces can be read as modes of internal gauge
exchange. In QED-like vocabulary an interpretation close to ``exchange
of virtual photons'' is possible, but this is not a particle
identification derived directly from the algebra of this paper.

\begin{center}\rule{0.5\linewidth}{0.5pt}\end{center}

\subsection{14. Hierarchy from the Flux Being
Complex}\label{hierarchy-from-the-flux-being-complex}

An important feature of this construction is that the face flux \(K_f\)
is itself a complex number:

\[
K_f\in\mathbb C.
\]

Hence the external flux of a condensate can be treated as a new
relational quantity at the next level.

Schematically,

\[
Z_e\in\mathbb C
\longrightarrow
K_f\in\mathbb C
\longrightarrow
K_{\mathcal C}\in\mathbb C
\longrightarrow\cdots
\]

That is, without changing the type

\[
\boxed{\mathbb C\to\mathbb C\to\mathbb C},
\]

we can read recursively

\[
\text{relations}\to\text{faces}\to\text{condensates}\to\text{relations between condensates}.
\]

This property allows small closed structures to be glued into larger
closed structures --- hierarchical condensation.

\begin{figure}
\centering
\pandocbounded{\includegraphics[keepaspectratio,alt={Fig. 4 Internal cancellation and the hierarchical growth of condensates}]{figs/fig4_internal_cancellation_condensate_en.png}}
\caption{Fig. 4 Internal cancellation and the hierarchical growth of
condensates}
\end{figure}

\emph{Fig. 4 Oppositely oriented contributions of shared relations
cancel internally, and a closed set of faces is read from outside as a
single condensate. Since the face flux \(K_f\in\mathbb C\) can be reused
as a relation at the next level, the hierarchy stays within the same
complex-number type.}

\begin{center}\rule{0.5\linewidth}{0.5pt}\end{center}

\subsection{15. Complete Closure and the Indistinguishability of
``Universe'' and ``Grandchild
Universe''}\label{complete-closure-and-the-indistinguishability-of-universe-and-grandchild-universe}

Suppose the relational flux between a closed condensate and its exterior
has vanished completely.

More strongly than a mere zero total: if no boundary relation readable
as an external connection exists, or if the entire effective flux on the
external boundary has vanished, the system is relationally cut off from
the outside.

For an internal observer, there is then no way to distinguish, without
external relations, whether the system is ``the whole universe'' or ``a
grandchild universe generated inside a larger system.''

Therefore the readout

\[
\boxed{\text{a completely closed relational world carries no absolute label of “the whole universe”}}
\]

becomes possible.

This is a physical hypothesis; it is not a derivation of actual
existence from the complex-number algebra alone.

\begin{center}\rule{0.5\linewidth}{0.5pt}\end{center}

\subsection{16. Composition from Past to Present and the Direction of
Time}\label{composition-from-past-to-present-and-the-direction-of-time}

The same complex-sum structure can also be used for a temporal readout.

Let the squared quantities fixed in the past be

\[
Z_1,\ldots,Z_n.
\]

Then the present composite quantity

\[
K_n=\sum_{j=1}^{n}Z_j
\]

is uniquely computable.

After the next relation \(Z_{n+1}\) is realized,

\[
K_{n+1}=K_n+Z_{n+1}.
\]

In this sense, no matter how many history items there are, the update
contracts to the two-body composition

\[
\boxed{(K_n,Z_{n+1})\to K_{n+1}}.
\]

On the other hand, there is no rule that uniquely factorizes or
generates the future \(Z_{n+1}\) from the present \(K_n\) alone.
Therefore,

\[
\boxed{\text{past}\to\text{present}}
\]

is definable as composition, whereas

\[
\boxed{\text{present}\to\text{future}}
\]

requires an additional generation rule.

This paper does not present this as a proof of the arrow of time itself,
but retains it as the \textbf{asymmetry between the composition
direction and the generation direction}.

\begin{center}\rule{0.5\linewidth}{0.5pt}\end{center}

\subsection{17. An Auxiliary Relation to the Log
Readout}\label{an-auxiliary-relation-to-the-log-readout}

For the multiplicative relation of nonzero complex numbers \(z_1,z_2\),

\[
z_3=z_1z_2,
\]

we have \(z_3\neq0\), and with properly managed branches the complex
logarithm gives

\[
\operatorname{Log} z_3=\operatorname{Log} z_1+\operatorname{Log} z_2.
\]

For the self-product

\[
z\mapsto z^2,
\]

take \(z\neq0\) and, on the chosen logarithmic branch, set

\[
q:=\operatorname{Log}z.
\]

Then in Log space

\[
q\mapsto 2q,
\]

which reads as additive self-doubling.

On the other hand, to read the sum of squared quantities

\[
\sum_n Z_n=K
\]

in Log space, set for each nonzero \(Z_n\)

\[
q_n:=\operatorname{Log}Z_n.
\]

Zero terms contribute nothing to the sum and can be excluded. Only when
\(K\neq0\),

\[
\operatorname{Log} K=\operatorname{Log}\left(\sum_n e^{q_n}\right).
\]

For \(K=0\), \(\operatorname{Log}K\) is undefined.

This Log readout is not essential to the loop/flux construction of this
paper, but it can be retained as an auxiliary readout connecting the
multiplicative foundation with additive observational representations.

\begin{center}\rule{0.5\linewidth}{0.5pt}\end{center}

\subsection{18. History of Deleted Hidden
Assumptions}\label{history-of-deleted-hidden-assumptions}

If only the conclusions of this paper were presented, it would look like
existing discrete gauge theory rewritten in complex numbers. We
therefore state explicitly the assumptions actually removed during the
thought experiment.

{\def\LTcaptype{none} % do not increment counter
\begin{longtable}[]{@{}
  >{\raggedright\arraybackslash}p{(\linewidth - 6\tabcolsep) * \real{0.2500}}
  >{\raggedright\arraybackslash}p{(\linewidth - 6\tabcolsep) * \real{0.2500}}
  >{\raggedright\arraybackslash}p{(\linewidth - 6\tabcolsep) * \real{0.2500}}
  >{\raggedright\arraybackslash}p{(\linewidth - 6\tabcolsep) * \real{0.2500}}@{}}
\toprule\noalign{}
\begin{minipage}[b]{\linewidth}\raggedright
Candidate
\end{minipage} & \begin{minipage}[b]{\linewidth}\raggedright
Initial reading
\end{minipage} & \begin{minipage}[b]{\linewidth}\raggedright
Verdict
\end{minipage} & \begin{minipage}[b]{\linewidth}\raggedright
Reason
\end{minipage} \\
\midrule\noalign{}
\endhead
\bottomrule\noalign{}
\endlastfoot
Complete network \(K_N\) & every relation connects to every relation &
deleted & not derivable from triangular loops alone; a 4-vertex, 5-edge
system with one diagonal also has loop solutions \\
Ratio \(z_2/z_1\) & minimal two-body relation & removed from the
foundation & easily smuggles in a comparison standard at readout \\
Equality test \(z_1=z_2\) & primitive same/different test & removed from
the foundation & no need to place an equality detector as a prior
physical operation \\
\(\sum z_n^2=0\) for two bodies & condition for two waves & deleted &
zero closure is not imposed as an extra condition in this thought
experiment \\
Triangle & a priori geometry of three bodies & deleted & readable as a
loop after \(W_1+W_2+W_3=0\) \\
Vertex sharing & connection principle of edges & on hold & vertices are
not presupposed; we moved to the dual edge-to-face readout \\
Adding a curvature variable & a field separate from the loop residual &
unnecessary & \(K_f=C_f^2\) reads as a complex central-projection
radius \\
Adding gauge particles & separate particles carrying flux & demoted to
physical interpretation & local flux propagation may be read as gauge
exchange, but independent particles are not derived \\
Completeness of many-body information & \(Z_e\) and \(K_f\) exhaust all
relational information of a many-body system & on hold & implicitly
assumed in v1.1; noticed to be underived on the occasion of Anegawa et
al.~{[}6{]}. It has not been shown that irreducible higher-order
relations do not exist independently \\
\end{longtable}
}

As the table shows, this paper moved in the direction of deriving or
deleting assumption candidates rather than adding assumptions.

\begin{figure}
\centering
\pandocbounded{\includegraphics[keepaspectratio,alt={Fig. 5 Method flowchart of the thought experiment}]{figs/fig5_thought_experiment_flowchart_en.png}}
\caption{Fig. 5 Method flowchart of the thought experiment}
\end{figure}

\emph{Fig. 5 The method of the thought experiment. Starting from the
closure of complex numbers, at each stage we check whether a complete
network, comparison by ratio, equality tests, zero closure, or
presupposed triangles have been assumed, shave off the unnecessary
assumptions, and adopt only the readouts that remain.}

\begin{center}\rule{0.5\linewidth}{0.5pt}\end{center}

\subsection{19. Axiom Candidates Remaining at This
Point}\label{axiom-candidates-remaining-at-this-point}

Apart from the algebra of complex numbers itself, the following
candidates remain at this stage.

\subsubsection{19.1 Simultaneous existence of multiple complex
relations}\label{simultaneous-existence-of-multiple-complex-relations}

It cannot yet be said that ``multiple occurrences'' have been generated
from a single value \(z\in\mathbb C\) by algebra alone. Therefore the
existence, or a generation rule, of multiple relations is an axiom
candidate.

\subsubsection{19.2 Generation of orientation and
incidence}\label{generation-of-orientation-and-incidence}

The structure in which a shared edge carries opposite signs in two faces
follows automatically once oriented boundaries are defined. But whether
the oriented incidence structure itself necessarily emerges from
nameless complex relations is unsettled.

\subsubsection{19.3 Reading a loop as a
face}\label{reading-a-loop-as-a-face}

Drawing the zero sum of three complex relations as a closed triangle is
mathematically natural, but identifying it with a physical ``face'' is a
readout rule.

\subsubsection{19.4 Reading a closed face complex as a
condensate}\label{reading-a-closed-face-complex-as-a-condensate}

Reading an internally cancelled closed polyhedron as one effective
object is the main physical interpretation candidate of this paper.

\subsubsection{19.5 Completeness of many-body information via pair
relations and loop
residuals}\label{completeness-of-many-body-information-via-pair-relations-and-loop-residuals}

Although not stated explicitly in v1.1, parts of the argument implicitly
used the assumption that

\[
\mathcal I_{\mathrm{pair+loop}}:=\{Z_e,K_f\}
\]

suffices to represent the relational information of a many-body system.
However, the completeness

\[
\boxed{\{Z_e,K_f\}\stackrel{?}{\Longrightarrow}\mathcal R_{\mathrm{full}}}
\]

is unproven. In particular, whether irreducible higher-order relational
quantities such as

\[
R_{123},\quad R_{1234},\quad\ldots
\]

are independently necessary, or whether they can be derived from \(Z_e\)
and \(K_f\), is unsettled. This paper does not adopt this point as a new
axiom but \textbf{keeps it on hold as an axiom candidate / derivation
problem.}

These are the objects of future study: whether they can be derived
further from complex-number algebra and namelessness.

\begin{center}\rule{0.5\linewidth}{0.5pt}\end{center}

\subsection{20. Connection to Existing Gauge
Theories}\label{connection-to-existing-gauge-theories}

Comparing with existing theories after obtaining the construction above,
the structural similarity with lattice gauge theory becomes clear.

In Wilson's lattice gauge theory, link variables on a discrete lattice
and closed loops play the central role {[}1{]}. In Hamiltonian lattice
gauge theory as well, electric flux and link structure are fundamental
{[}2{]}.

This paper does not put a gauge group or a lattice first, but reading

\[
K_f=\sum_{e\in\partial f}\epsilon_{fe}Z_e
\]

as a face flux, the format ``quantities on edges → loops → residuals on
faces'' is close to commutative lattice gauge theory.

The difference lies in the starting point. In this paper:

\begin{enumerate}
\def\labelenumi{\arabic{enumi}.}
\tightlist
\item
  No lattice is assumed.
\item
  No complete graph is assumed.
\item
  No vertices are assumed first.
\item
  No specific gauge group is assumed first.
\item
  Both squared quantities and face residuals belong to the same
  \(\mathbb C\).
\end{enumerate}

Hence this is not a simplified construction of existing gauge theory;
the order is that gauge-theoretic structure was connected only after
following the closure of complex numbers.

\begin{center}\rule{0.5\linewidth}{0.5pt}\end{center}

\subsection{21. Comparison with Closure Defects, the Bianchi Identity,
Bubble Networks, and Multi-partite
Entanglement}\label{comparison-with-closure-defects-the-bianchi-identity-bubble-networks-and-multi-partite-entanglement}

Livine relaxes the closure constraint in the coarse-graining of spin
networks and interprets the closure defect as a local measure of
curvature inside the coarse-grained region {[}3{]}.

Charles and Livine reinterpret the closure constraint of the hyperbolic
tetrahedron as a Bianchi identity, connecting face holonomies with
polyhedral closure {[}4{]}.

The bubble networks of Freidel and Livine glue three-dimensional
geometric cells along interfaces and handle the structure of boundary
faces, fluxes, and gluing {[}5{]}.

These have strong similarity to this paper's

\[
\text{relations}\to\text{loops}\to\text{residuals}\to\text{face flux}\to\text{closed cells}.
\]

The differences are equally clear. In the prior work, structures such as
spin networks, \(SU(2)\) holonomies, spinors, graphs, and cell complexes
exist from the start. This paper does not take them as the starting
point; it examines how far one can reach the same type of readout from

\[
z\in\mathbb C
\]

and

\[
Z=z^2\in\mathbb C
\]

together with the closure of complex addition.

The relation to prior work is therefore positioned not as ``imitation of
the same structure'' but as \textbf{verification of the possibility of
having reached the same kind of structure independently from fewer
premises}.

Anegawa, Miyata, Suzuki, Tamaoka {[}6{]} showed that in a fixed
tripartition of the two-sided BTZ black hole, a quantity measuring
genuine multi-partite entanglement keeps growing even after the ordinary
entanglement entropies of the subsystems saturate, continuing to track
the growth of the Einstein--Rosen bridge.

This result cannot be identified with the \(K_f\) or the condensates of
this paper, nor does it directly support or refute the present
complex-relation model. What matters for this paper is that it supplies
\textbf{a concrete example that low-order readouts need not exhaust the
total information of a many-body system}. Through it we recognized the
implicit assumption remaining in v1.1 --- ``\(Z_e\) and \(K_f\) suffice
to exhaust the effective relational information of a many-body system''
--- and in v1.2 returned that completeness to the list of open problems.

\begin{center}\rule{0.5\linewidth}{0.5pt}\end{center}

\subsection{22. Candidate Novelty}\label{candidate-novelty}

We divide the candidate novelty of this paper into mathematical novelty
and interpretive novelty.

\subsubsection{22.1 Mathematical novelty}\label{mathematical-novelty}

No novelty is claimed for the closure of complex numbers, sums of
squares, closed polygons, or boundary cancellation themselves.

\subsubsection{22.2 Interpretive novelty
candidates}\label{interpretive-novelty-candidates}

What this paper proposes is the following readout chain:

\[
\boxed{
\mathbb C
\to
z^2\in\mathbb C
\to
\sum Z_n=K\in\mathbb C
\to
\text{closure / defect}
\to
\text{face flux}
\to
\text{closed condensate}
\to
\text{hierarchical condensate}
}
\]

Moreover, in obtaining this chain we assumed no complete network, no
fixed vertices, no fixed dimension, and no external comparison standard.

\subsubsection{22.3 Methodological novelty
candidate}\label{methodological-novelty-candidate}

In this paper, instead of adding axioms while building a model, we
delete a hidden assumption each time one is discovered.

What matters in the end is therefore

\[
\boxed{\text{not what was introduced, but what never had to be introduced.}}
\]

\begin{center}\rule{0.5\linewidth}{0.5pt}\end{center}

\subsection{23. Limitations and Open
Problems}\label{limitations-and-open-problems}

This paper alone does not derive the following.

\begin{enumerate}
\def\labelenumi{\arabic{enumi}.}
\tightlist
\item
  How multiple independent complex relations are generated from a single
  complex number.
\item
  Whether the oriented incidence structure necessarily arises from
  namelessness.
\item
  Which loops are selected as actually existing faces.
\item
  Whether triangular loops suffice, or higher polygonal loops are
  independently necessary.
\item
  The time-evolution law of the flux.
\item
  How flux propagation maps to the individual gauge particles of the
  Standard Model.
\item
  Stability conditions of condensates and the readout of mass, energy,
  and spin.
\item
  Whether completely closed systems correspond to actually existing
  independent universes.
\item
  The geometric meaning of complex \(K_f\), in particular of its
  imaginary part, and the correspondence from the complex
  central-projection radius to general non-constant curvature fields.
\item
  Whether the set of pairwise relations \(Z_e\) and loop residuals
  \(K_f\) suffices to determine the full relational information of a
  many-body system. This informational sufficiency can be tested by
  fixing all face residuals and the external boundary readout on the
  face--edge incidence matrix and computing, by a rank computation, the
  dimension of the solution space admitted for the internal \(Z_e\).
\item
  Whether irreducible relational invariants of three or more bodies,
  \(R_{123},R_{1234},\ldots\), exist independently, underivable from
  pairwise and loop quantities.
\end{enumerate}

Items 10 and 11 in particular make explicit, in v1.2, the hidden
assumption that remained in v1.1. Their validity is currently unsettled
and is left to future theoretical analysis and numerical experiments.

These are not negative evidence but connection problems for the next
stage.

\begin{center}\rule{0.5\linewidth}{0.5pt}\end{center}

\subsection{24. Conclusion}\label{conclusion}

The starting point of this paper is one ordinary complex number:

\[
z\in\mathbb C.
\]

A complex number squared is still a complex number:

\[
Z_n:=z_n^2\in\mathbb C.
\]

Therefore

\[
\sum_n z_n^2=C^2
\]

is, once squared quantities are read as basic,

\[
\boxed{\sum_n Z_n=K},
\]

an ordinary sum of complex numbers.

From this simple rereading, one reads a zero sum of three relations as a
loop, a nonzero sum as a loop residual, and by gluing oriented loops
obtains the cancellation of shared internal relations and the external
boundary flux. Since the flux is again complex, it can be reused as a
relation at the next level. A closed set of faces can be read as one
effective object with its internal degrees of freedom cancelled --- a
condensate.

What matters is that in obtaining these we did not first posit complete
networks, fixed vertices, triangles, gauge groups, or external scales.
Rather, by doubting them as assumption candidates and deleting those
that could not be derived, a loop/flux structure similar to existing
gauge theory appeared afterwards.

The central claim of this paper is therefore not a new equation:

\[
\boxed{
\text{Give the trivial closure of complex numbers a relational readout with minimized assumptions,}
\text{and the hierarchy of loops, residuals, fluxes, and condensates can be read in a unified way.}
}
\]

However, this unified readout does \textbf{not claim to exhaust the full
relational information of many-body systems.} In v1.2 we withdrew the
implicit assumption that the pairwise relations \(Z_e\) and the loop
residuals \(K_f\) suffice, and left as an open problem whether
irreducible higher-order relations are independently necessary.

We do not claim that this framework is a final physical theory. But as a
minimal basis for clarifying which additional assumptions are truly
indispensable, it merits examination.

\begin{center}\rule{0.5\linewidth}{0.5pt}\end{center}

\subsection{References}\label{references}

This paper uses no self-citation. The following are external references
consulted, after the construction of this paper was obtained, to check
structural similarities with existing theories.

\begin{enumerate}
\def\labelenumi{\arabic{enumi}.}
\tightlist
\item
  K. G. Wilson, ``Confinement of quarks,'' \emph{Physical Review D}
  \textbf{10}, 2445--2459 (1974). DOI: 10.1103/PhysRevD.10.2445.
\item
  J. Kogut and L. Susskind, ``Hamiltonian formulation of Wilson's
  lattice gauge theories,'' \emph{Physical Review D} \textbf{11},
  395--408 (1975). DOI: 10.1103/PhysRevD.11.395.
\item
  E. R. Livine, ``Deformation Operators of Spin Networks and
  Coarse-Graining,'' \emph{Classical and Quantum Gravity} \textbf{31},
  075004 (2014). arXiv:1310.3362. DOI: 10.1088/0264-9381/31/7/075004.
\item
  C. Charles and E. R. Livine, ``The closure constraint for the
  hyperbolic tetrahedron as a Bianchi identity,'' \emph{General
  Relativity and Gravitation} \textbf{49}, 92 (2017). arXiv:1607.08359.
  DOI: 10.1007/s10714-017-2255-2.
\item
  L. Freidel and E. R. Livine, ``Bubble Networks: Framed Discrete
  Geometry for Quantum Gravity,'' \emph{General Relativity and
  Gravitation} \textbf{51}, 9 (2019). arXiv:1810.09364. DOI:
  10.1007/s10714-018-2493-y.
\item
  T. Anegawa, A. Miyata, S. Suzuki, and K. Tamaoka, ``Growing
  Einstein--Rosen Bridge with Multi-partite Entanglement,''
  arXiv:2609.18225 (2026).
\end{enumerate}

\begin{center}\rule{0.5\linewidth}{0.5pt}\end{center}

\subsection{Appendix A: The Minimal Series of
Formulas}\label{appendix-a-the-minimal-series-of-formulas}

Extracting only the algebraic skeleton of this paper gives the following
series.

\[
z_n\in\mathbb C
\]

\[
Z_n:=z_n^2\in\mathbb C
\]

\[
K:=\sum_n Z_n\in\mathbb C
\]

\[
K=C^2,\qquad C\in\mathbb C
\]

Loop:

\[
\sum_{e\in\partial f}\epsilon_{fe}Z_e=0
\]

Loop residual:

\[
\sum_{e\in\partial f}\epsilon_{fe}Z_e=K_f
\]

Central-projection-radius readout:

\[
K_f=C_f^2
\]

Shared-edge cancellation:

\[
(+Z_e)+(-Z_e)=0
\]

Composition of multiple faces:

\[
\sum_f K_f=\sum_{e\in\partial\Omega}\epsilon_e Z_e
\]

Global identity on a closed oriented face complex (derived identity
structure):

\[
\boxed{\sum_{f\in F(\mathcal C)}K_f=0}
\]

\(F(\mathcal C)\) is the set of oriented faces constituting the
condensate \(\mathcal C\). On a closed polyhedral boundary in which
every edge belongs to two faces with opposite orientations, this holds
identically by incidence cancellation, i.e., by the discrete
\(\partial^2=0\).

Hierarchization:

\[
K_f\in\mathbb C\quad\Rightarrow\quad K_f\ \text{is reusable as a relational quantity at the next level}
\]

\begin{center}\rule{0.5\linewidth}{0.5pt}\end{center}

\subsection{Appendix B: Ledger of Axiom
Candidates}\label{appendix-b-ledger-of-axiom-candidates}

{\def\LTcaptype{none} % do not increment counter
\begin{longtable}[]{@{}
  >{\raggedright\arraybackslash}p{(\linewidth - 4\tabcolsep) * \real{0.3333}}
  >{\raggedright\arraybackslash}p{(\linewidth - 4\tabcolsep) * \real{0.3333}}
  >{\raggedright\arraybackslash}p{(\linewidth - 4\tabcolsep) * \real{0.3333}}@{}}
\toprule\noalign{}
\begin{minipage}[b]{\linewidth}\raggedright
Object
\end{minipage} & \begin{minipage}[b]{\linewidth}\raggedright
Current classification
\end{minipage} & \begin{minipage}[b]{\linewidth}\raggedright
Reason
\end{minipage} \\
\midrule\noalign{}
\endhead
\bottomrule\noalign{}
\endlastfoot
\(z\in\mathbb C\) & mathematical background & complex numbers adopted as
the descriptive language \\
\(z^2\in\mathbb C\) & derived identity structure & closure of complex
numbers \\
\(\sum z_n^2=C^2\) & derived identity structure &
\(\sum z_n^2\in\mathbb C\) and existence of square roots \\
Reading \(Z_n=z_n^2\) as the basic quantity & readout definition & entry
point of the new physical interpretation \\
Existence of multiple relations & axiom candidate & multiple occurrences
not yet derived from a single complex number \\
Reading a loop as a face & readout candidate & natural as a geometric
representation, but physical identity unproven \\
Oriented incidence & axiom candidate / derivation problem & needed for
shared-edge cancellation; derivation from namelessness incomplete \\
Reading \(K_f=C_f^2\) as a complex central-projection radius & geometric
readout & re-expression of the same sum-of-squares structure \\
Closed polyhedron = condensate & physical interpretation candidate & has
internal cancellation and external effectivization \\
Local flux = gauge exchange & connection problem with existing theory &
similar to lattice gauge theory, but particle identity underived \\
Completeness of many-body information
\(\{Z_e,K_f\}\Rightarrow\mathcal R_{\mathrm{full}}\) & axiom candidate /
derivation problem & implicitly assumed in v1.1 but unproven; the
possibility of independently existing irreducible higher-order
relational quantities cannot be excluded \\
Complete network & deleted & not derivable from loop conditions alone \\
\end{longtable}
}

\end{document}
